%% 中文简历模板

\documentclass[12pt,a4paper,utf8]{moderncv}
%\usepackage{CJKutf8}

\usepackage{ctex}
\CTEXnoindent

% moderncv themes
%\moderncvtheme[green]{casual}                 % optional argument are 'blue' (default), 'orange', 'red', 'green', 'grey' and 'roman' (for roman fonts, instead of sans serif fonts)
\moderncvtheme[green]{classic}                % idem
%
%% character encoding
%\usepackage[utf8]{inputenc}                   % replace by the encoding you are using

% adjust the page margins
\usepackage[scale=0.8]{geometry}
%\setlength{\hintscolumnwidth}{3cm}						% if you want to change the width of the column with the dates
%\AtBeginDocument{\setlength{\maketitlenamewidth}{6cm}}  % only for the classic theme, if you want to change the width of your name placeholder (to leave more space for your address details
%\AtBeginDocument{\recomputelengths}                     % required when changes are made to page layout lengths

% personal data
\firstname{{~}\\ An-Bao Xu  \huge{（Ph.D.) Associate Professor} }
\familyname{}
\title{}  %
%\title{Resumé title (optional)}               % optional, remove the line if not wanted
%\address{500 Dongchuan Rd}{200241 Shanghai}    % optional, remove the line if not wanted
\address{Wenzhou University}{College of Mathematics and Physics, Office：3A405}
%\address{办公室：3A405}      %

%\mobile{15057702302}                    % optional, remove the line if not wanted
%\phone{15057702302}                      % optional, remove the line if not wanted
%\fax{0577-86689098}                          % optional, remove the line if not wanted
\email{xuanbao@wzu.edu.cn}                      % optional, remove the line if not wanted
%\homepage{math.wzu.edu.cn/Art/Art_190/Art_190_1648.aspx}                % optional, remove the line if not wanted
%\extrainfo{additional information (optional)} % optional, remove the line if not wanted
\photo[64pt]{photo.jpg}                         % '64pt' is the height the picture must be resized to and 'picture' is the name of the picture file; optional, remove the line if not wanted
%\quote{\textbf{多一点奉献, 社会不至于退步}}                 % optional, remove the line if not wanted
%\quote{\textbf{多一点奉献, 社会更美好！}}                 % optional, remove the line if not wanted

% to show numerical labels in the bibliography; only useful if you make citations in your resume
\makeatletter
\renewcommand*{\bibliographyitemlabel}{\@biblabel{\arabic{enumiv}}}
\makeatother

% bibliography with mutiple entries
%\usepackage{multibib}
%\newcites{book,misc}{{Books},{Others}}

%\nopagenumbers{}                             % uncomment to suppress automatic page numbering for CVs longer than one page
%----------------------------------------------------------------------------------
%            content
%----------------------------------------------------------------------------------
\begin{document}
\CTEXnoindent
\renewcommand{\baselinestretch}{1.0}
\maketitle

%\section{基本介绍}
%\cvline{}{1985年出生于浙江瑞安，现为温州大学新湖青年学者，特聘教授，硕士生导师，浙江省中青年学科带头人。
%研究领域为应用统计，一直致力于寿命数据分析、医学数据分析、贝叶斯计算以及客观贝叶斯方法等领域的研究。
%已主持完成国家自然科学基金青年项目一项，浙江省自然科学基金面上项目一项，中国博士后自然科学基金及陕西省博士后科学基金各一项，
%目前正在主持国家自然科学基金面上项目一项。}

\section{Education}
\cventry{2013.09-2017.07}{Ph.D.}{Hunan University,  College of Mathematics}{}{}{}%{City}{\textit{Grade}}{Description}  % arguments 3 to 6 can be left
\cventry{2010.09-2013.07}{M.S.}{Guilin University of Electronic Technology,  School of Mathematics and Computing Science}{}{}{}
\cventry{2006.09-2010.07}{B.S.}{Xianning College（Hubei University of Science and Technology）}{}{}{}
%\section{博士论文}
%\cvline{论文题目}{\emph{屏蔽数据的贝叶斯统计分析}}
%\cvline{导师}{汤银才 教授}
%\cvline{研究内容}{从工程实际背景出发, 提出了用贝叶斯方法来对屏蔽数据进行建模, 放宽已有文献中
%的三个基本假设：1、系统中的各个部件相互独立, 且寿命分布已知; 2、产品的寿命试验在使用应力下进行的;
%3、所考虑的产品是串联系统. 首次在可靠性统计中提出用客观贝叶斯方法来估计模型中的参数,
%并考察了估计的频率性质; 首次提出用非参数贝叶斯方法来估计系统的生存函数，放宽了独立性以及寿命分布已知的假设;
%首次考虑了在恒加寿命试验以及步加寿命试验下屏蔽数据的贝叶斯统计分析. 本文还研究了屏蔽概率假设的检验问题以
%及并联系统中屏蔽数据的贝叶斯统计分析.}
%
%本文所提供的步加与序加试验统计分析方法,不仅具有重要的理论意义,而且
%上述方法在固体钽电解电容器上的成功应用说明它们
%具有更深远的现实意义:它使加速寿命试验的统计分析方法更加成熟,
%为步加试验与序加试验在产品可靠性分析中的推广与应用提供理论上
%的保证,并将给可靠性工程人员提供方法论上的指导.}

\section{Professional Information}
\subsection{Work Experience}
%\cventry{2017.7-至今}{瓯江特聘教授}{温州大学应用统计与金融计量研究所}{}{}{}
%\cventry{2013.12-至今}{副教授}{温州大学应用统计与金融计量研究所}{}{}{}
\cventry{2023.12.-present}{Associate Professor}{Wenzhou University, \ College of Mathematics and Physics}{}{}{}
\cventry{2017.07-2023.12}{Lecturer
}{Wenzhou University, \ College of Mathematics and Physics}{}{}{}

\subsection{Academic Visit}
\cventry{2016.11-2017.05}{Visiting Student}{Chinese Academy of Sciences}{}{}{}
\cventry{2014.08-2016.08}{Joint PhD candidate}{Auburn University, USA}{}{}{}
%\cventry{2010.9-2011.3}{访问学者}{美国密苏里大学哥伦比亚分校}{}{}{}

%\subsection{教学经历}
%\cventry{2017.07-至今}{讲授课程}{}{}{}{}
%\cvline{ }{$\circ$高等数学 \qquad \qquad \qquad \qquad \qquad \qquad \qquad $\circ$高等代数}
%\cvline{ }{$\circ$实用计算方法（数值分析） \qquad \qquad \qquad \ $\circ$数学与人类文明}

\section{Research Interests}
%\subsection{工作经历}
\cvline{ 1}{Numerical Linear Algebra}
\cvline{ 2}{Privacy Computing}
\cvline{ 3}{Optimization}
\cvline{ 4}{Deep Learning}
\cvline{ 5}{Scientific Computing}
%\cvline{ 4}{温州大学第四届“学生科技创新优秀指导教师”(2014)}
%\cvline{ 5}{温州大学第四届“我心目中的好导师”(2015)}
%\cvline{6}{首届中国运筹学会可靠性分会优秀论文奖（2017）}
%\cvline{ 7}{浙江省高校中青年学科带头人 (2017)}

\section{ Participation and Hosting of Foundation}
%\cventry{2017.1--2020.12}{竞争失效模型中若干关键问题的研究}{国家自然科学基金（面上项目）}{}{直接经费48万}{\textcolor{blue}{主持}}
\cventry{2019.01--2021.12}{Research on the Parameterized Thresholding Algorithm for Compressed Sensing and Matrix Completion}{National Natural Science Foundation of China（Youth Program）}{(11801418)}{}{\textcolor{blue}{Host}}
\cventry{2014.08--2016.08}{Joint PhD candidate}{National Scholarship for Building High Level University}{(201406130046)}{}{\textcolor{blue}{Host}}
%\cventry{2015.9--2016.9}{可修屏蔽系统可靠性的贝叶斯评估方法}{中国博士后科学基金面上项目}{(2015M572598)}{5 万}{\textcolor{blue}{主持}}
%\cventry{2015.9--2016.9}{屏蔽数据系统可靠性评估的贝叶斯方法}{陕西省博士后科学基金一等资助}{}{4 万}{\textcolor{blue}{主持}}

%\section{横向课题}
%\cventry{2018.9--2018.11}{基于贝叶斯方法的隐患挖掘预警算法技术}{华为终端技术合作部}{}{经费总计21.63万（与华东师范大学合作）}{}

%\cventry{2018.8--2018.10}{基于贝叶斯方法的隐患挖掘预警算法技术合作（华为终端技术合作部）}{}{}{\textcolor{blue}{到账经费10万}}

%\section{企事业培训}
%\cventry{2018.8}{可靠性数据分析}{蔚来汽车可靠性部门}{}{16课时}{}
%\cventry{2017.4}{医学数据分析}{宁波卫计委}{}{8课时}{}


%\section{研究兴趣}
%\cvline{}{贝叶斯计算、竞争风险模型、退化数据分析、寿命数据分析、可靠性}
%\cvline{\textbf{I.}}{\textbf{寿命数据分析:医学数据以及产品寿命数据。}提出时空模型来分析温州手足口病数据，为温州市疾控中心
%提出非参数贝叶斯法以及多层贝叶斯法有效地解决失效机理屏蔽时系统可靠性的估计。}
%\cvline{\textbf{II.}}{\textbf{竞争风险模型：}提出Frailty双维纳退化过程，解决系统可靠度以及剩余寿命可靠度没显示表达的问题；提出多层贝叶斯方法识别维纳退化过程的变点问题；提出基于最优线性无偏估计准则改进几类退化过程的在线估计。}
%\cvline{\textbf{III.}}{\textbf{寿命试验最优设计：}基于Kullback-Liebler距离，从信息论角度提出新的贝叶斯最优设计准则，并证明基于该准则得到的最优设计方案具有全局最优性，同时提出大样本近似法以及蒙特卡洛法来寻找基于该最优准则的最优设计方案，这是国际上首次从信息论角度研究加速寿命试验最优设计。}



\section{Journal Articles（$\ast$ Corresponding Author)：}
\begin{enumerate}

\item  Jun Lei, Jiqian Zhao, Jingqi Wang, \textbf{An-Bao Xu}*, \emph{Tensor completion via leverage sampling and tensor QR decomposition for network latency estimation}, Applied Intelligence, 55:687(2025), 16p.

\item Frank Uhlig, \textbf{An-Bao Xu}*,  \emph{Iterative optimal solutions of linear matrix equations for Hyperspectral and Multispectral image fusing}, Calcolo, 60:26(2023), 33 p.

\item Yongming Zheng, \textbf{An-Bao Xu}*,  \emph{Tensor completion via tensor QR decomposition and $L_{2,1}$-norm minimization}, Signal Processing, 189(2021), 10 p. 
\item Dongxiu Xie, Hugo J. Woerdeman, \textbf{An-Bao Xu}*,  \emph{Parametrized quasi-soft thresholding operator for compressed sensing and matrix completion}, Computational and Applied Mathematics, 149(2020), 24 p.
\item \textbf{An-Bao Xu}*, Dongxiu Xie,  \emph{Low-rank approximation pursuit for matrix completion}, Mechanical Systems and Signal Processing, 95(2017), 77-89. %（SCI二区）
\item Dongxiu Xie, \textbf{An-Bao Xu}*, Zhen-yun Peng,  \emph{Least-squares symmetric solution to the matrix equation AXB=C with the norm inequality constraint}, International Journal of Computer Mathematics, 93(2016), 1564-1578.%（SCI一区）
\item \textbf{An-Bao Xu}, Zhen-yun Peng*,  \emph{Norm-constrained least-squares solutions to the matrix equation AXB = C}, Abstract and Applied Analysis, 2013 (2013), 10 p.

\item Jiahao Guo, \textbf{An-Bao Xu}*,  \emph{High-performance privacy-preserving matrix completion for trajectory recovery}, arXiv preprint arXiv: 	arXiv:2405.05789 (2024).

\item Jiqian Zhao, \textbf{An-Bao Xu}*,  \emph{Infrared small target detection via tensor $L_{2,1}$-norm minimization and ASTTV regularization: a novel tensor recovery approach}, arXiv preprint arXiv:2402.18003 (2024).



\item Xiaohui Ni, \textbf{An-Bao Xu}*,  \emph{A randomized algorithm for a QR decomposition-based approximate SVD}, arXiv preprint arXiv: arXiv:2305.11450  (2023).



%\item Jie Li, Yanjun Fu, \textbf{Ancha Xu}$^\ast$, Zumu Zhou and Weiming Wang, A Spatial-Temporal ARMA Model of the Incidence of Hand, Foot, and Mouth Disease in Wenzhou, China, Abstract and Applied Analysis, Vol. 2014, Article ID 238724.（SCI二区）
%\item Qiang Guang,Yincai Tang and \textbf{Ancha Xu}$^\ast$, Objective Bayesian Analysis For Bivariate Marshall-Olkin Exponential Distribution, Computational Statistics and Data Analysis, Vol. 64(7), pp.299-313, 2013.（SCI二区）
%\item \textbf{Ancha Xu}$^\ast$ and Yincai Tang, Objective Bayesian Analysis of Accelerated Competing Failure Models under Type-I Censoring. Computational Statistics and Data Analysis, Vol. 55(10), pp.2830-2839, 2011.（SCI二区）
%\item \textbf{Ancha Xu}$^\ast$ and Yincai Tang, Bayesian analysis of Birnbaum-Saunders distribution with partial information, Computational Statistics and Data Analysis, Vol. 55(7), pp.2324-2333, 2011.（SCI二区）
%\item \textbf{Ancha Xu} and Yincai Tang$^\ast$, Reference Analysis for Birnbaum-Saunders Distribution, Computational Statistics and Data Analysis, Vol. 54(1), pp.185-192, 2010.（SCI二区）
%\item \textbf{Ancha Xu}$^\ast$, Lijuan Shen ,Bingxing Wang, and Yincai Tang, On modeling bivariate Wiener degradation process, IEEE Transactions on Reliability, Vol. 67(3), pp. 897-906, 2018. （可靠性工程顶级期刊）
%\item Pingping Wang, Yincai Tang, Suk Joo Bae and \textbf{Ancha Xu}$^\ast$, Bayesian approach for two-phase degradation data based on change-point wiener process with measurement errors,  IEEE Transactions on Reliability, Vol. 67(2), pp. 688-700, 2018.（可靠性工程顶级期刊）
%\item Piao Chen, \textbf{Ancha Xu}$^\ast$ and Zhisheng Ye, Generalized fiducial inference for accelerated life tests with Weibull distribution and progressively type-II censoring, IEEE Transactions on Reliability. Vol. 65(4), pp.1737-1744, 2016.（可靠性工程顶级期刊）
%\item \textbf{Ancha Xu}$^\ast$ and Yincai Tang, A Bayesian method for planning accelerated life testing, IEEE Transactions on Reliability. Vol. 64(4), pp.1383-1392, 2015.（可靠性工程顶级期刊）
%\item \textbf{Ancha Xu}$^\ast$, Sanjib Basu and Yincai Tang, A full Bayesian approach for masked data in step-stress accelerated life testing, IEEE Transactions on Reliability, Vol. 63(3), pp.798-806, 2014.（可靠性工程顶级期刊）
%\item \textbf{Ancha Xu} and Yincai Tang$^\ast$, Bayesian analysis of Pareto reliability with dependent masked data, IEEE Transactions on Reliability, Vol. 58(4), pp.583-588, 2009.（可靠性工程顶级期刊）
%\item \textbf{Ancha Xu}, Yincai Tang, Qiang Guan and Xinze Lian$^\ast$, Planning simple step-stress accelerated life tests using reference optimality criterion, Proceedings of the iMeche, Part O: Journal of Risk and Reliability. Vol. 230(1), pp.85-92, 2016.（SCI三区）
%\item \textbf{Ancha Xu}$^\ast$ and Yincai Tang, Reference optimality criterion for planning accelerated life testing, Journal of Statistical Planning and Inference, Vol. 167, pp.14-26, 2015.（SCI三区）
%\item Shengliang Lu, Jie Zhang, Shirong Zhou and \textbf{Ancha Xu} $^\ast$, 
%Reliability prediction of the axial ultimate bearing capacity of piles: A hierarchical Bayesian method, Advances in Mechanical Engineering, Vol.10(11), pp. 1-11, 2018. （SCI四区）
%\item \textbf{Ancha Xu} and Lijuan Shen$^\ast$, Improved on-line estimation for gamma process, Statistics and Probability Letters,
%Vol.143, pp. 67-73, 2018. （SCI四区）
%\item Yan Shen and \textbf{Ancha Xu}$^\ast$, On the dependent competing risks using Marshall-Olkin bivariate Weibull model: parameter estimation with different methods, Communications in Statistics – Theory and Methods. Vol. 47(22), pp. 5558-5572, 2018. （SCI四区）
%\item \textbf{Ancha Xu}$^\ast$ and Shirong Zhou, Bayesian analysis of series system with dependent causes of failure, Statistical Theory and Related Fields. Vol. 1(1), pp.128-140, 2017. (创 刊号)
%\item Yalan Li and \textbf{Ancha Xu}$^\ast$, Fiducial inference for Birnbaum-Saunders distribution, Journal of Statistical Computation and Simulation. Vol. 86(9), pp.1673-1685, 2016.（SCI四区）
%\item \textbf{Ancha Xu}$^\ast$, Yincai Tang and Qiang Guan, Bayesian Analysis of Masked Data in Step- stress Accelerated Life Testing, Communications in Statistics - Simulation and Computation, Vol. 43(8), pp.2016-2030, 2014.（SCI四区）
%\item Qiang Guan, Yincai Tang, Jiayu Fu and \textbf{Ancha Xu}$^\ast$, Optimal Multiple Constant-stress Accelerated Life Tests For Generalized Exponential Distribution, Communications in Statistics - Simulation and Computation, Vol. 43(8), pp.1852-1865, 2014.（SCI四区）
%\item \textbf{Ancha Xu}$^\ast$ and Yincai Tang, Posterior Propriety in Nonparametric Mixed Effects Model, Applied Mathematics---A Journal of Chinese Universities, Series B, Vol.28(3), pp.369-378, 2013. (SCI四区）
%\item Jiayu Fu, \textbf{Ancha Xu}$^\ast$ and Yincai Tang, Objective Bayesian Analysis of Pareto Distribution under Progressive Type-II Censoring, Statistics and Probability Letters, Vol. 82(10), pp.1829-1836, 2012.（SCI四区）
%\item \textbf{Ancha Xu} and Yincai Tang$^\ast$, Objective Bayesian Analysis for Linear Degradation Models, Communications in Statistics – Theory and Methods, Vol.41(21), pp.4034-4046, 2012.（SCI四区）
%\item \textbf{Ancha Xu} and Yincai Tang$^\ast$, Nonparametric Bayesian analysis of competing risks problem with masked data, Communications in Statistics – Theory and Methods, Vol. 40(13), pp.2326-2336, 2011.（SCI四区）
\end{enumerate}

%\section{Invited Talks：}
%\begin{enumerate}
%	\item The 8th International Conference on Matrix Analysis and Application (ICMAA 2019), University of Nevada, Reno, USA, July 15-18, 2019. \\
%	-Contributed Talk: \emph{Tensor completion via low-rank approximation pursuit}.
%	
%	\item 2018 SIAM Conference on Applied Linear Algebra (SIAM-ALA18), Hong Kong Baptist University, Hong Kong, China, May 4-8, 2018. \\
%	-Contributed Talk: \emph{Parametrized quasi-soft thresholding operator for compressed sensing and matrix completion}.
%	\item 2017 Workshop on Matrices and Operators (MAO), Hunan University, Changsha, Hunan, China, June 9-12, 2017.
% \\
%	-Contributed Talk: \emph{Varied Parametric Quasi-Soft Thresholding for Compressed Sensing}.
%	
%	\item 2016 SIAM Conference on Imaging Science (SIAMIS16), Albuquerque, New Mexico, USA, May 23-26, 2016.
%\\
%	-Contributed Talk: \emph{Low-Rank Approximation Pursuit for Matrix and Tensor Completion}.
%\\
%	-Travel Support from SIAM Student Travel Fund.
%	
%	\item The 40th SIAM Southeastern Atlantic Section Conference (SIAM-SEAS), University of Georgia, Athens, Georgia, USA, March 12-13, 2016.
%\\
%	-Contributed Talk: \emph{Optimal Solutions for Solvable and Unsolvable Inhomogeneous Linear Matrix Equations of Generalized Sylvester Type}.
%	
%\end{enumerate}

%\section{科研学会会员}
%\cvline{1}{中国数学会}
%\cvline{2}{中国工业与应用数学学会 (CSIAM)}
%\cvline{3}{Society for Industrial and Applied Mathematics (SIAM)}
%\cvline{4}{International Linear Algebra Society (ILAS)}
	

%\section{指导硕士生}
%\cvline{2019级}{郑永铭}
%\cvline{2012,2014-2016年}{美国大学生数学建模竞赛二等奖}
%\cvline{2012年}{全国大学生数学建模竞赛浙江省一等奖}
%\cvline{2011,2014,2017年}{全国研究生数学建模竞赛三等奖}

\end{document}


%% end of file `template_en.tex'.
